\documentclass[reqno]{amsart} 
\usepackage{amssymb,latexsym,amsmath,amscd,graphicx,setspace,amsthm,verbatim,enumitem}
\usepackage[margin = 3 cm]{geometry}

\def\upint{\mathchoice%
    {\mkern13mu\overline{\vphantom{\intop}\mkern7mu}\mkern-20mu}%
    {\mkern7mu\overline{\vphantom{\intop}\mkern7mu}\mkern-14mu}%
    {\mkern7mu\overline{\vphantom{\intop}\mkern7mu}\mkern-14mu}%
    {\mkern7mu\overline{\vphantom{\intop}\mkern7mu}\mkern-14mu}%
  \int}
\def\lowint{\mkern3mu\underline{\vphantom{\intop}\mkern7mu}\mkern-10mu\int}


\theoremstyle{plain}
\newtheorem{theorem}{Theorem}[section]
\newtheorem{proposition}{Proposition}
\newtheorem{corollary}{Corollary}
\newtheorem{lemma}{Lemma}
\newtheorem{conjecture}{Conjecture}
\newtheorem{question}{Question}
\newtheorem{problem}{Problem}
      
\theoremstyle{definition}
\newtheorem{definition}{Definition}

\newenvironment{solution1}{\paragraph{\emph{Solution $1$}.}}{\hfill$\square$}
\newenvironment{solution2}{\paragraph{\emph{Solution $2$}.}}{\hfill$\square$}
\newenvironment{solution3}{\paragraph{\emph{Solution $3$}.}}{\hfill$\square$}

\begin{document} 

\title[Homework 7]{Homework 7}

\date{\today} 
\maketitle 


Let $(X,\mathcal{A},\mu)$ be a measure space for which $\mu(X) < \infty$ and let $f:X \rightarrow \mathbb{R}$ be a bounded function.  In class we defined what it means for $f$ to be Lebesgue integrable in which case we also defined the meaning of the Lebesgue integral
$$\int_{X}f. $$
The Lebesgue satisfies several properties listed below.  You should try to prove as many of those properties as possible.
\begin{problem}
Assume $f, g \in B(X)$ are both Lebesgue integrable.  Then $f+g$ is Lebesgue integrable and
$$\int_{X}(f+g) = \int_{X}f + \int_{X}g. $$
\end{problem}
\begin{solution1}

\end{solution1}

\begin{problem}
Assume $f \in B(X)$ is Lebesgue integrable and $\lambda \in \mathbb{R}$.  Then $\lambda f$ is also Lebesgue integrable and
$$\int_{X}(\lambda f) = \lambda \int_{X}f. $$
\end{problem}
\begin{solution1}

\end{solution1}


\begin{problem}
If $f,g \in B(X)$ are Lebesgue integrable and $f \le g$, then 
$$\int_{X}f \le \int_{X}g. $$
Moreover, $|f|$ is Lebesgue integrable as well and
$$\Big|\int_{X}f \Big| \le \int_{X} |f|. $$
\end{problem}
\begin{solution1}

\end{solution1}


\begin{problem}
If $f \in B(X)$ is Lebesgue integrable and $A, B \in \mathcal{A}$ are such that 
$$X = A \cup B \text{ and } \mu(A \cap B) = 0, $$
then $f$ is Lebesgue integrable on both $A$ and $B$ and
$$\int_{X}f = \int_{A}f + \int_{B}f. $$
\end{problem}
\begin{solution1}

\end{solution1}








\end{document}









